%% feature/sec-methods
\section{Method}
\label{sec:method}

\subsection{Overview}
\label{sec:overview}

We isolate a single variable: how many times a fixed, already-useful insight is
repeated before a problem. Related work leaves the sign of this effect genuinely
open (\S\ref{subsec:rw-positioning})---repetition helps non-reasoning models but is
near-null with reasoning enabled, the activation-redistribution mechanism predicts
help for content-light repetition yet harm for diverse semantic insertions, and raw
reasoning traces actively degrade accuracy as more are added. We therefore state two
directional hypotheses and test them under a design in which the repetition count
$k$ is the only quantity that changes.

\begin{description}
  \item[\textbf{H1 (Accuracy).}] Repeating the insight more than once
        ($k>1$) increases task accuracy over the single-insertion baseline
        ($k{=}1$), as repeated exposure reinforces the already-useful guidance.
  \item[\textbf{H2 (Cost).}] Repetition decreases total generated tokens, and
        hence cost, because the reinforced insight displaces some of the model's
        own re-derivation; the added input tokens are processed in the
        parallelizable prefill stage rather than in generation
        \citep{leviathan2025prompt, zhao2026trs}.
\end{description}

Because the prompt-repetition efficiency claim is that repetition shifts cost to
prefill rather than generation, and because distilled insights are reported to act
primarily on generated tokens \citep{zhao2026trs}, we treat token efficiency as a
first-class axis alongside accuracy rather than a secondary diagnostic. We report
both jointly: accuracy gains bought at a disproportionate token cost, or token
savings that sacrifice accuracy, each falsify the combined H1--H2 picture even if one
metric improves in isolation. The remainder of this section makes the intervention,
data, models, and accounting precise.

% NOTE: the old \subsection{Research Question} is removed here -- it restated the
% abstract and the question is already captured by the Overview + H1/H2. The
% one-sentence framing is folded into the Overview's first paragraph.

\subsection{Task Formulation}
\label{sec:task-formulation}

Let $\mathcal{D} = \{(q_i, a_i, e_i)\}_{i=1}^{n}$ be a benchmark dataset, where $q_i$ is a mathematical problem, $a_i$ is the gold final answer, and $e_i$ is a precomputed insight associated with the problem. For each target item $i$, we construct a prompt
\[
    x_i^{(k)} = \operatorname{Prompt}\bigl(q_i, \underbrace{e_i, \ldots, e_i}_{k}\bigr),
\]
where $k \in \mathcal{K} = \{0, 1, 2, 3, 5\}$ is the repetition count. Here $k=0$ is the no-insight baseline, $k=1$ is the standard single-insight setting, and $k>1$ tests whether repeating the same insight adds benefit or causes degradation.

Given a language model $M_\theta$, the model produces $\hat{y}_i^{(k)} = M_\theta(x_i^{(k)})$, and we score correctness against the gold answer,
\[
    c_i^{(k)} = \mathbb{1}\bigl[\operatorname{Eval}(\hat{y}_i^{(k)}, a_i) = 1\bigr],
\]
yielding the accuracy at repetition count $k$,
\[
    \operatorname{Acc}(k) = \frac{1}{n}\sum_{i=1}^{n} c_i^{(k)}.
\]


\subsection{Repetition Intervention and Prompt Variants}
\label{sec:repetition-intervention}
% Preamble: \usepackage{tikz}  \usetikzlibrary{arrows.meta}
\begin{figure*}[t]
\centering
\resizebox{\textwidth}{!}{%
\begin{tikzpicture}[
    font=\small,
    >={Latex[length=2mm]},
    box/.style={draw=#1!70!black, fill=#1!12, rounded corners=2pt,
                align=center, inner sep=5pt, minimum height=11mm},
    card/.style={draw=violet!70!black, fill=violet!10, rounded corners=2pt,
                 align=center, inner sep=2pt, minimum width=40mm,
                 minimum height=5.5mm, font=\footnotesize},
    ctx/.style={card, draw=black!45, fill=black!6},
    verb/.style={draw=orange!70!black, fill=orange!12, rounded corners=2pt,
                 align=center, inner sep=1pt, minimum width=40mm,
                 minimum height=4mm, font=\scriptsize\itshape},
    brace/.style={decorate, decoration={brace, amplitude=4pt}, thick}
]

% ============================================================
% Variant (a): insight-only repetition
% ============================================================
\begin{scope}[shift={(0,0)}]
  \draw[draw=violet!55, dashed, rounded corners=3pt]
      (-2.55, 2.30) rectangle (2.55, -2.55);
  \node[font=\bfseries] at (0, 2.65)
      {(a) Insight-only};
  \node[font=\footnotesize, text=black!60] at (0, 2.05)
      {context, question, insight$\times k$};

  \node[ctx] (a-sys)  at (0, 1.40) {System prompt};
  \node[ctx] (a-q)    at (0, 0.75) {Problem $q_i$};
  \node[card] (a-i1)  at (0, 0.05) {Insight $e_i$};
  \node[card] (a-i2)  at (0,-0.55) {Insight $e_i$};
  \node[card] (a-i3)  at (0,-1.40) {Insight $e_i$};
  \node[font=\footnotesize] at (0,-0.95) {$\vdots$};

  % brace marking the repeated block
  \draw[brace] (2.05, 0.30) -- (2.05, -1.65)
      node[midway, right=5pt, font=\scriptsize] {$k\times$};
\end{scope}

% ============================================================
% Variant (b): full-block repetition
% ============================================================
\begin{scope}[shift={(6.6,0)}]
  \draw[draw=violet!55, dashed, rounded corners=3pt]
      (-2.55, 2.30) rectangle (2.55, -2.55);
  \node[font=\bfseries] at (0, 2.65)
      {(b) Full-block};
  \node[font=\footnotesize, text=black!60] at (0, 2.05)
      {(context, question, insight)$\times k$};

  % block 1
  \draw[draw=violet!40, rounded corners=2pt] (-2.15, 1.62) rectangle (2.15, 0.18);
  \node[ctx, minimum width=34mm] (b-q1) at (0, 1.30) {Problem $q_i$};
  \node[card, minimum width=34mm] (b-i1) at (0, 0.55) {Insight $e_i$};
  % block 2
  \draw[draw=violet!40, rounded corners=2pt] (-2.15, -0.05) rectangle (2.15, -1.49);
  \node[ctx, minimum width=34mm] (b-q2) at (0, -0.37) {Problem $q_i$};
  \node[card, minimum width=34mm] (b-i2) at (0, -1.12) {Insight $e_i$};
  % continuation
  \node[font=\footnotesize] at (0,-1.85) {$\vdots$};

  \draw[brace] (2.30, 1.62) -- (2.30, -1.49)
      node[midway, right=5pt, font=\scriptsize] {$k\times$};
\end{scope}

% ============================================================
% Variant (c): full-block repetition, verbose
% ============================================================
\begin{scope}[shift={(13.2,0)}]
  \draw[draw=violet!55, dashed, rounded corners=3pt]
      (-2.55, 2.30) rectangle (2.55, -2.55);
  \node[font=\bfseries] at (0, 2.65)
      {(c) Full-block, verbose};
  \node[font=\footnotesize, text=black!60] at (0, 2.05)
      {(b) $+$ connecting phrase};

  % block 1
  \draw[draw=violet!40, rounded corners=2pt] (-2.15, 1.62) rectangle (2.15, 0.18);
  \node[ctx, minimum width=34mm] (c-q1) at (0, 1.30) {Problem $q_i$};
  \node[card, minimum width=34mm] (c-i1) at (0, 0.55) {Insight $e_i$};
  % verbose connector
  \node[verb] (c-v) at (0, -0.18) {``Let me repeat that''};
  % block 2
  \draw[draw=violet!40, rounded corners=2pt] (-2.15, -0.52) rectangle (2.15, -1.85);
  \node[ctx, minimum width=34mm] (c-q2) at (0, -0.84) {Problem $q_i$};
  \node[card, minimum width=34mm] (c-i2) at (0, -1.53) {Insight $e_i$};
  \node[font=\footnotesize] at (0,-2.18) {$\vdots$};

  \draw[brace] (2.30, 1.62) -- (2.30, -1.85)
      node[midway, right=5pt, font=\scriptsize] {$k\times$};
\end{scope}

% ============================================================
% Shared footnote
% ============================================================
\node[font=\footnotesize, text=black!60] at (6.6, -3.15)
    {Insight $e_i$ is the same fixed skill card in every copy; only $k$ and the
     variant change. Evaluated at $k \in \{1,2,3,4,5\}$.};

\end{tikzpicture}}
\caption{Three repetition variants for the distilled insight. \textbf{(a)}
Insight-only: the context and question appear once and only the skill card $e_i$
is repeated $k$ times, isolating re-exposure of the insight. \textbf{(b)}
Full-block: the entire (context, question, insight) block is repeated $k$ times,
mirroring the canonical prompt-repetition design but with the insight included.
\textbf{(c)} Full-block verbose: as in (b), with a connecting phrase introducing
each repeated block, following the verbose prompt-repetition variant. In all
variants the insight text, question, gold answer, model, and decoding settings are
held fixed.}
\label{fig:variants}
\end{figure*}

The intervention is deliberately minimal: for each problem we take one fixed insight
text $e_i$ and copy it $k$ times; no new insight is generated when $k$ changes. Any
difference between conditions is therefore caused by repeated exposure to the same
text, not by a different insight, a better extraction, or any change in content.

How the insight is repeated may matter as much as how many times, so we evaluate
three variants that differ in what is held constant across copies
(Figure~\ref{fig:variants}). Let the prompt consist of a context, the question, and
the insight.

\begin{enumerate}
  \item[\textbf{(a)}] \textbf{Insight-only repetition:} context and question appear
        once, and only the insight block is repeated $k$ times
        (context, question, $\text{insight} \times k$). This is the tightest test of
        the repetition effect, isolating re-exposure of the insight from any change
        to the rest of the prompt.
  \item[\textbf{(b)}] \textbf{Full-block repetition:} the entire
        (context, question, insight) block is repeated $k$ times, mirroring the
        canonical prompt-repetition design of \citet{leviathan2025prompt} but with
        the insight included in each copy.
  \item[\textbf{(c)}] \textbf{Full-block repetition, verbose:} as in (b), but each
        repeated block is introduced with a connecting phrase
        (e.g., ``Let me repeat that''), following the verbose variant of
        \citet{leviathan2025prompt}, which sometimes outperformed plain repetition.
\end{enumerate}

Across all variants we hold the question, gold answer, insight text, prompt template,
model, and decoding settings fixed, so $k$ and the variant are the only quantities
that change. We evaluate $k \in \{0,1,2,3,5\}$: $k{=}0$ is the no-insight baseline,
$k{=}1$ the single insertion, $k{=}2$ matches the strongest setting reported in
prompt-repetition work, $k{=}3$ tests whether any effect persists beyond that
optimum, and $k{=}5$ probes over-repetition and prompt-bloat failure modes. The exact
templates for all variants and baselines are given in
Appendix~\ref{app:prompt-templates}; variant~(a) is the template shown below.


\begin{quote}
\ttfamily\footnotesize
You are solving a math problem.\\[2pt]
Useful insight:\\
<skill\_text>\\[2pt]
Useful insight:\\
<skill\_text>\\
\ldots\\[2pt]
Problem:\\
<question>\\[2pt]
Return the final answer in boxed form.
\end{quote}

For $k=0$ the prompt contains only the instruction and the problem; for $k>0$ the
same \texttt{skill\_text} appears exactly $k$ times.


\subsection{Dataset}
\label{sec:dataset}

We use the released TRS dataset, \texttt{stallone0000/Rea\-soning-Skill} on Hugging
Face,\footnote{\url{https://huggingface.co/datasets/stallone0000/Reasoning-Skill}}
linked to the paper \emph{Thinking with Reasoning Skills: Fewer Tokens, More
Accuracy}.\footnote{\url{https://arxiv.org/abs/2604.21764}} Each row already provides
exactly the three ingredients our study requires---a target problem, a reference
answer, and a precomputed skill card:
\[
    (q_i, a_i, e_i) = (\texttt{question}, \texttt{answer}, \texttt{skill\_text}).
\]

Reusing this dataset is the core methodological choice, because it removes the main
alternative explanation for any result we might observe. Building insights ourselves
would require a demonstration pool, an embedding index, a retrieval policy, and an
extraction prompt, and any effect of repetition could then be confounded with
retrieval or extraction quality. Since the skill cards already exist and are
independently validated to help at $k=1$, we reuse them exactly as provided and treat
each card as a fixed intervention unit. The only quantity that changes across
conditions is the repetition count $k$, rather than the compound of
\text{retriever} + \text{extractor} + \text{quality} + k.

We focus first on the DeepMath/TRS math slice exposed through the release, evaluating
on a non-overlapping split where one is available. AIME and GSM8K can be added later
as robustness checks, but each requires compatible precomputed skill cards, so we
treat them as secondary.

\subsection{Models and Decoding}
\label{sec:model-selection}

We evaluate several reasoning-capable models through OpenRouter, selected on three
criteria: the model is exposed as a reasoning model, the API returns token-usage
metadata, and where possible it reports reasoning-token usage separately from visible
output. The families we consider are Qwen thinking variants, DeepSeek reasoning models
(e.g.\ R1), and Kimi/Moonshot thinking models, with OpenAI-compatible reasoning models
as additional reference points where their accounting permits. Spanning more than one
family lets us separate model-specific behavior from effects that generalize.

For every model, decoding settings are held fixed across repetition counts:
deterministic decoding at temperature $0$, a fixed maximum output length, and a fixed
reasoning-effort setting (e.g.\ \texttt{medium}). This ensures that differences
between $k$ values are not driven by sampling variation or a shifting reasoning
budget.

\subsection{Token and Cost Accounting}
\label{sec:token-accounting}

For each response we log input tokens $t^{(k)}_{i,\mathrm{in}}$ and total completion
tokens $t^{(k)}_{i,\mathrm{out}}$, with total tokens
\[
    t^{(k)}_{i,\mathrm{total}} = t^{(k)}_{i,\mathrm{in}} + t^{(k)}_{i,\mathrm{out}}.
\]
For reasoning models, completion tokens already include any explicit or hidden
reasoning tokens; when the provider reports them separately we additionally log
$t^{(k)}_{i,\mathrm{reason}}$ as a breakdown of $t_{\mathrm{out}}$, not as an additive
term. For example, $t_{\mathrm{out}}=1600$ with $t_{\mathrm{reason}}=1300$ means the
reasoning tokens are contained within the $1600$, and hence within
$t_{\mathrm{total}}$.

Because repetition inflates the input while the insight is meant to compress the
output, total tokens are our primary efficiency axis. We report accuracy together with
token, and where available cost, efficiency:
\begin{align*}
    \operatorname{Eff}_{\mathrm{tok}}(k)  &= \frac{\operatorname{Acc}(k)}{\operatorname{MeanTotalTokens}(k)}, \\[4pt]
    \operatorname{Eff}_{\mathrm{cost}}(k) &= \frac{\operatorname{Acc}(k)}{\operatorname{MeanCost}(k)}.
\end{align*}

\subsection{Evaluation Protocol}
\label{sec:evaluation}

We use deterministic, answer-based evaluation rather than an LLM judge. Each output is
postprocessed to extract the final answer, preferably from a
\texttt{\textbackslash boxed\{...\}} expression, and the normalized prediction is
compared against the normalized gold answer. For every item and repetition count we
store the question id, the gold and predicted answers (raw and normalized),
correctness, and the full token and cost breakdown.

The design is fully paired: each problem is evaluated under every repetition count, so
we can analyze per-item transitions in addition to aggregate accuracy. For two counts
$k_a$ and $k_b$ we tabulate how many items move between outcomes,
\[
    N_{\mathrm{w}\to\mathrm{r}},\ \,
    N_{\mathrm{r}\to\mathrm{w}},\ \,
    N_{\mathrm{r}\to\mathrm{r}},\ \,
    N_{\mathrm{w}\to\mathrm{w}},
\]
each as a function of $(k_a, k_b)$. This matters because aggregate accuracy can hide
whether repetition helps the same problems consistently, trades some correct answers
for others, or only moves a few borderline cases.

All conditions are run with $3$ independent random seeds (controlling prompt-order
shuffling and any stochastic decoding). We report mean accuracy and mean token counts
with standard deviations. Significance between conditions at matched $k$ is assessed
with a paired McNemar test \citep{mcnemar1947} at $p<0.1$, following
\citet{leviathan2025prompt}.
